\chapter*{摘\quad 要}
\addcontentsline{toc}{chapter}{摘要}

大语言模型 Agent 已从文本生成扩展到文件读写、Shell 执行、网络访问和持久 Memory 等环境交互。相应地，安全问题不再局限于输入或输出内容，而必须覆盖真实副作用的产生、暂存、提交、取消与恢复。\Aegis{} 面向这一需求实现了事务式安全执行 Runtime：所有工具调用经过风险分类、策略与权限判定、受控执行、检查、提交或回滚，并将全过程记录为可追溯的 Audit 事件。

本作品的三项核心创新为：（1）\textbf{Transactional Controlled Execution}：将可恢复副作用置于 Pending 状态，并以 Checkpoint、Commit Gate 和 Effect 生命周期约束可信状态的形成；（2）\textbf{Approval-aware TaskGraph Scheduling}：在任务图中仅阻断等待审批节点的依赖或资源冲突范围，使安全独立节点能够继续推进；（3）\textbf{Dependency-aware Selective Rollback}：回滚受影响依赖闭包内的 Effect，保留未进入受影响范围且已安全提交的独立 Effect。Risk、Policy、Permission、Approval 与 Audit 是上述运行时链路的支撑机制，而非对核心创新的替代命名。

最终验收覆盖静态检查、后端、前端和真实浏览器流程：Ruff 通过，Pyright 为 0 errors、0 warnings；后端测试为 889 passed、8 skipped；前端 Vitest 为 14 passed；Playwright 为 11 passed。正式实验数据固定为 Safety 225 行、Graph 18 行和 Rollback 105 行，并由原始数据及聚合结果重建核验。Safety 使用确定性受控规则，不以 LLM 裁决作为正式结论依据；Graph 在固定时延、确定性受控环境中验证工程行为，不作统计显著性主张。

\vspace{1em}
\noindent\textbf{关键词：} Agent 安全；事务式受控执行；审批感知调度；选择性回滚；状态完整性

\chapter{作品概述}

\section{作品背景}

Agent 将语言模型的决策接入外部工具后，系统可能读取不可信文档、修改文件、执行命令或形成跨任务的持久状态。间接提示注入已经表明，第三方内容可能被模型错误解释为高优先级指令，从而诱导集成应用执行非预期操作\cite{greshake2023indirect}。代码 Agent 的风险评测也将危险命令和执行环境置于问题中心\cite{guo2024redcode}。因此，安全 Runtime 必须把控制点放在副作用发生的位置，而不是仅依赖对输入、输出或模型意图的单点判断。

\section{问题定义与威胁模型}

本作品处理的对象是携带工具调用请求和任务契约的 Agent 运行时。对于每一项动作，系统需要判断它是否符合用户目标与授权边界、是否允许在安全结论确定前进入隔离状态，以及失败或取消时应恢复哪些状态。为此，系统从以下四个层面描述威胁模型：

\begin{itemize}
  \item \textbf{指令层：}网页、文档与工具输出中的不可信内容可能携带间接提示注入或语义操纵；
  \item \textbf{资源层：}路径越界、危险 Shell、敏感文件访问与未授权外发可能直接产生高风险副作用；
  \item \textbf{状态层：}未经验证的工具结果或 Memory 写入若直接成为可信状态，可能造成跨任务持续污染\cite{dash2026memorypoisoning}；
  \item \textbf{调度层：}审批等待、取消与恢复若未正确约束，可能造成兄弟节点泄漏、孤儿执行或重复执行\cite{khan2026stop}。
\end{itemize}

系统将 Planner 输出视为不可信提案。可信边界由 Runtime Scheduler、风险策略与权限判定、审批门、Effect 生命周期和 Audit 记录共同构成；工具原始输出不直接成为可信 Memory，也不以原始不可信内容展示给前端。

\section{国内外研究现状}

研究界已从风险发现逐步走向 Agent 运行时治理。ToolEmu 通过模拟工具环境识别工具使用中的长尾风险\cite{ruan2024toolemu}；R-Judge 与 Agent Security Bench（ASB）分别聚焦交互轨迹风险意识和多场景攻防评测\cite{yuan2024rjudge,zhang2025asb}；AgentHarm 则从多步恶意任务角度评估 Agent 被滥用的风险\cite{andriushchenko2025agentharm}。这些工作表明，安全评估需要覆盖多步轨迹和真实工具影响，而不能只观察文本拒答。

在运行时防护方面，SafeAgent 将安全建模为随轨迹演化的有状态决策问题\cite{liu2026safeagent}；SafeHarness 将控制机制置于 Agent 生命周期和执行 Harness 中\cite{lin2026safeharness}；Consent Integrity 强调审批内容必须由可信中介绑定到实际执行对象\cite{weng2026consent}。Cordon 已将工具意图、结果血缘、暂存 Effect 与审计元数据置于任务级语义事务之中\cite{chen2026cordon}。这些工作说明，\Aegis{} 不将“事务”或“审批”本身宣称为首创，而将重点放在三项机制在可运行 Runtime 中的受控组合、工程闭环与可复现验证。

\section{设计目标}

\Aegis{} 的设计目标是使安全判断、任务推进和状态恢复处于同一条可验证链路中。首先，所有真实副作用经由统一 Runtime Scheduler 进入受控执行，风险未确定或权限不足时采用 Fail-Closed 行为。其次，可恢复的文件、下载与 Memory 状态遵循 Pending 到可信状态的生命周期；检查失败、拒绝或取消时，系统据已冻结的关联关系恢复受影响范围。再次，审批不成为全局暂停信号：仅依赖等待节点或与其 Effect target 冲突的节点被阻断，无依赖且无冲突的节点可继续调度。最后，Risk、Policy、Permission、Approval、Checkpoint、Effect、Commit/Rollback 和 Audit 都形成结构化事实，以支持全过程审计可追溯性。

\section{适用范围与边界}

\Aegis{} 适用于需要工具调用、运行时审计、人工审批和受控恢复的本地或企业自动化工作流，例如安全敏感代码辅助、文件处理和多步骤业务准备流程。当前原型未覆盖分布式调度、服务重启恢复、完整 ACID 事务、真实支付或邮件发送、通用 MCP/Skill 协议，以及同一节点 speculative execution。对外发类动作，系统以受控、脱敏和干运行边界为前提；本作品不据此声称对开放网络环境中的任意副作用完成通用回滚。
